\chapter{AI Prompt Log}
\ChapterLabel{AIPromptLog}
\section{Chapter Responsibility}
\label{Section:AIPromptLogResponsibility}
Preserve supplied facts, AI drafting scope, and human-review boundaries.

\noindent\textbf{Expected result:} A reviewable project record with explicit evidence and gaps.


\section{AIS-02 Initialization: November 9, 2026}
The user asked Codex to instantiate template 2.33 in a separate directory,
review the facts and gaps before dependent edits, keep separate TikZ sources
and trace links, and return the complete build for review. The user confirmed
the project name, namespace, mission, scope, system boundary, stakeholders,
owners, and Embedded Systems applicability. The user specified the
publication destination and reserved acceptance and publication authorization
for a later step.

The confirmed facts and the prior project draft informed the candidate
chapters. Candidate requirements, priorities, and test designs remain draft.
Technical and evidence decisions are tracked in
\hyperref[Table:OwnedQuestions]{the owned questions}.

\section{Context Draft and Drawing (November 14, 2026)}
The user requested a context description and matching drawing that identify
external devices, the host application, the consumer, and what crosses each
boundary. The \hyperref[Section:MeasurementMission]{context draft} distinguishes
the control app boundary from the system boundary and links the missing
engineering results to the \hyperref[Table:OwnedQuestions]{owned questions}.

\section{Conversation Logging (November 19, 2026)}
After recovering three archived conversations and exporting timestamped
Markdown transcripts, the user requested an automatic project rule for
future conversations. \hyperref[Rule:AIChatLogging]{The conversation
transcript rule} defines the destination, verbatim content, timestamp format,
refresh points, and active-response limitation. The local exporter refreshes
available project conversations without committing or publishing them.

\section{Integrated Specification Publication (November 23, 2026)}
The user authorized restructuring the publication into a System Requirements
and Architecture Specification, retaining complete requirements and detailed
test, verification, validation, traceability, and evidence registers.
The publication scope is defined in
\path|s8sDocumentGuide.tex|. Existing artifact IDs and obligations are
preserved; unresolved decisions and Pending evidence remain explicit.

\section{Project Co-authorship (November 23, 2026)}
The user requested Chadwick Adebiyi as co-author of all project documents.
\hyperref[Section:ProjectAuthorship]{Project Authorship} records the decision;
the shared \texttt{DocumentAuthors} metadata supplies publication and
transcript-header credits. Repository authoring guidance applies the same
credit to future project documents while retaining verbatim historical
message bodies and original external-source attribution.

\section{Recording and Playback Stories (November 23, 2026)}
\label{Section:RecordingPlaybackPrompt}
The user asked which stakeholder stories could support recording sensor data
and playback for debugging. After reviewing the proposed stories, the user
accepted stories 5, 6, 7, 8, and 11 and requested their addition to the User
Stories chapter. AI skill \hyperref[Skill:AIS-05]{\texttt{AIS-05}} added
\UserStoryRef{US-05}, \UserStoryRef{US-06}, \UserStoryRef{US-07},
\UserStoryRef{US-08}, and \UserStoryRef{US-09} with their acceptance intent.
One agent performed the edit and self-review. No stories for result comparison
or recording export were added. No derived requirement approval, product
verification, stakeholder validation result, or release authorization is inferred.
The conversation's working directory is outside the source repository;
automatic transcript discovery does not capture this chat without an explicit
source export under \hyperref[Rule:AIChatLogging]{the transcript rule}.

The user subsequently requested consecutive story numbering and suggestions
for use cases and diagrams. The replay-output story originally proposed as
US-11 is now \UserStoryRef{US-09}; linked references were updated.
Use-case and diagram suggestions remain proposals pending selection.

\section{Recording and Playback Engineering Package (November 23, 2026)}
\label{Section:DiagnosticPackagePrompt}
The user requested requirements derived from the proposed use cases, requirement verifications and use-case validations, then authorized proceeding. One aiaCoordinator/aiaDeveloper performed the connected edits and self-review using \hyperref[Skill:AIS-01]{\texttt{AIS-01}}, \hyperref[Skill:AIS-05]{\texttt{AIS-05}}, \hyperref[Skill:AIS-06]{\texttt{AIS-06}}, \hyperref[Skill:AIS-07]{\texttt{AIS-07}}, \hyperref[Skill:AIS-08]{\texttt{AIS-08}}, \hyperref[Skill:AIS-09]{\texttt{AIS-09}}, \hyperref[Skill:AIS-11]{\texttt{AIS-11}}, \hyperref[Skill:AIS-13]{\texttt{AIS-13}} and \hyperref[Skill:AIS-14]{\texttt{AIS-14}}. No delegated or independent review is claimed. The plan was a sequential 15--25 minute authoring/checking effort: scenarios and diagrams approximately 5--10 minutes, 3,000--5,000 tokens; requirements and reciprocal plans 5--10 minutes, 4,000--7,000 tokens; compilation and self-review 5 minutes, 2,000--4,000 tokens, with uncertainty from layout and compiler diagnostics. Shared registers made delegation overhead unjustified. Quantitative limits, recording/compatibility policies, replay time/state and routing remain owned questions; all new requirements are Proposed and all execution dispositions Pending. Human technical review and stakeholder execution remain future work.

\section{Playback Review Corrections (November 25, 2026)}
\label{Section:PlaybackReviewCorrectionsPrompt}
The user requested review of PR 6, then authorized merging the changes into origin/main and updating the Architecture Specification. The review identified missing initialization on playback start/restart and ambiguous speed scaling. One aiaCoordinator/aiaDeveloper applied \hyperref[Skill:AIS-06]{\texttt{AIS-06}}, \hyperref[Skill:AIS-07]{\texttt{AIS-07}}, \hyperref[Skill:AIS-08]{\texttt{AIS-08}}, \hyperref[Skill:AIS-09]{\texttt{AIS-09}}, \hyperref[Skill:AIS-13]{\texttt{AIS-13}} and \hyperref[Skill:AIS-14]{\texttt{AIS-14}} sequentially, with self-review. No independent review, product execution, stakeholder acceptance or requirement-baseline approval is claimed. Initialization mechanisms and engineering policies remain owned decisions. The local transcript is refreshed under \hyperref[Rule:AIChatLogging]{the transcript rule}.

\section{Ordered Main Integration (November 25, 2026)}
The user clarified the workflow test: merge the Architecture Specification work first, then the updated user stories. The first merge reached origin/main as \texttt{e593772}. The user-story integration preserves the reader-publication structure, repository-only guidance, co-author metadata and generated \path|s8sArchitecture.pdf|. Conflicting history, prompt, evidence and scope text were reconciled without dropping either change. The superseded \path|s8sManual.pdf| remains removed. GitHub connector and CLI API merge attempts returned service errors before either merge; ordinary Git merges and pushes preserve the requested sequence.
